%----------------------------------------------------------------------------------------
%	PACKAGES AND OTHER DOCUMENT CONFIGURATIONS
%----------------------------------------------------------------------------------------
\DocumentMetadata{
pdfversion=2.0,  
lang=en-US,   
pdfstandard={ua-2, a-4f},
tagging = on, 
tagging-setup={math/setup={mathml-SE}} ,
}
\documentclass[12pt]{article}
\usepackage{amsmath}
\usepackage[extreme]{savetrees}
\usepackage[utf8]{inputenc}
\usepackage{newpxtext}
\usepackage{hyperref}
\usepackage{multicol}
\usepackage{physics}
\usepackage{siunitx}

%%%%%%%%% === Document Configuration === %%%%%%%%%%%%%%

\hypersetup{
pdftitle={MATH 340 -- Homework 10 Spring 2026},
pdfauthor={Ignacio Rojas},
pdfkeywords={tenth homework},%
}

\newcommand{\twobyone}[2]{\begin{pmatrix} %% 2 x 1 matrix
   #1 \\ #2 \end{pmatrix}}
\newcommand{\twobytwo}[4]{\begin{pmatrix} %% 2 x 2 matrix
   #1 & #2 \\ #3 & #4 \end{pmatrix}}
\newcommand{\threebyone}[3]{\begin{pmatrix} %% 3 x 1 matrix
   #1 \\ #2 \\ #3 \end{pmatrix}}
\newcommand{\threebythree}[9]{\begin{pmatrix} %% 3 x 3 matrix
   #1 & #2 & #3 \\ #4 & #5 & #6 \\ #7 & #8 & #9 \end{pmatrix}}
\newcommand{\cL}{\mathcal{L}}

%----------------------------------------------------------------------------------------
%	ARTICLE CONTENTS
%----------------------------------------------------------------------------------------
\begin{document}

\begin{center}
   {\Large MATH 340 -- Homework 10,\quad Spring 2026}
\end{center}

Recall that you must hand in a subset of the problems for which deleting any problem makes the total score less than 10. The maximum possible score on this homework is 10 points. See the syllabus for scoring details. NOTE: All problems from Boyce, DiPrima, and Meade are from the \(11^{\text{th}}\) edition, which is available online: \href{https://www.math.colostate.edu/~liu/MATH340_S26Crd/Textbook_BoyceDiPrimaMeade_11ed.pdf}{https://www.math.colostate.edu/\(\sim\)liu/MATH340\_S26Crd/Textbook\_BoyceDiPrimaMeade\_11ed.pdf}.

\begin{enumerate}

   \item (1) [1 point] For \(p>-1\) show, by definition of the Laplace transform, that

         \[
            \cL\lbrace t^p\rbrace(s)=\frac{\Gamma(p+1)}{s^{p+1}},
         \]

         where \(\Gamma(x)=\int_{0}^{\infty}e^{-t}t^{x-1}\dd t\).

   \item (1) [1 point] For \(a\) a constant, show, by definition of the Laplace transform, that

         \[
            \cL\lbrace e^{at}\rbrace(s)=\frac{1}{s-a}.
         \]

   \item (1+) [2 points] Find the Laplace transform of \(e^t\) by integrating term-wise the series
         \[
            e^{t}=\sum_{n=0}^{\infty}\frac{t^n}{n!}.
         \]

   \item (1+) [2 points] The trigonometric functions can be expressed as follows:
         \[
            \sin(x)=\frac{e^{ix}-e^{-ix}}{2i},\quad\cos(x)=\frac{e^{ix}+e^{-ix}}{2}.
         \]
         Find the Laplace transform of \(\sin\) and \(\cos\) via linearity, and by definition of the Laplace transform.

   \item (1+) [2 points] Find the Laplace transform of \(\sin\) and \(\cos\) by separating the Laplace transform of \(e^{it}\) into real and imaginary parts and applying Euler's identity.

   \item (1+) [2 points] Via integration by parts, verify that
         \[
            \cL\lbrace t^pe^{at}\rbrace=\frac{\Gamma(p+1)}{(s-a)^{p+1}}.
         \]

   \item (1+) [2 points] Find the Laplace transforms of the functions
         \[
            e^{at}\sin(bt),\quad\text{and}\quad e^{at}\cos(bt),
         \]
         where \(a,b\) are constants.

   \item (2) [3 points] Find the Laplace transform of \(\cos(2\sqrt{t})/\sqrt{\pi t}\) by expanding cosine into its Maclaurin series and integrating term-wise. Recall
         \[
            \cos(t)=\sum_{n=0}^{\infty}(-1)^n\frac{t^{2n}}{(2n)!}.
         \]
   \item (2+) [4 points] Let \(F(s)=\cL\lbrace f(t)\rbrace(s)\) be the Laplace transform of \(f(t)\). Show that
         \[
            \cL\lbrace e^{at}f(t)\rbrace=F(s-a),
         \]
         for a constant \(a\).

   \item (2+) [4 points] Let \(F(s)=\cL\lbrace f(t)\rbrace(s)\) be the Laplace transform of \(f(t)\). Show that
         \[
            \cL\lbrace f(at)\rbrace=\frac{1}{a}F\left(\frac{s}{a}\right),
         \]
         for a constant \(a\).

   \item (3) [9 points] Boyce, DiPrima, and Meade chapter 6 section 1 problem 23.

   \item (2+) [4 points] Define the \emph{beta function} as
         \[
            \beta(x,y)=\int_{0}^{1}t^{x-1}(1-t)^{y-1}\dd t.
         \]
         Show \(\beta\) is a \emph{symmetric} function, this is \(\beta(x,y)=\beta(y,x)\). Verify the identity \(\beta(x,y)=\frac{\Gamma(x)\Gamma(y)}{\Gamma(x+y)}\).

   \item (3-) [8 points] Research the definition of the Laplace transform for matrices of functions. Show that for a matrix \(A\) with complex entries,
         \[
            \cL\lbrace e^{At}\rbrace =\left(sI-A\right)^{-1}
         \]
\end{enumerate}

\end{document}